\section{Related Work}
\label{sec:relatedwork}

Our research intersects with three main areas of literature: Large Language Model agents in software engineering, program verification/validation pipelines, and distributed coordination and audit frameworks.

\subsection{LLM Agents in Software Engineering}
Since the advent of scalable Transformers~\cite{vaswani2017attention} and few-shot reasoning models like GPT-3~\cite{brown2020language}, code generation has moved from simple autocompletion to agentic orchestration. Frameworks such as Aider, AutoGPT, and Devin-like implementations have shown that wrapping LLMs in read-write-execute tool loops allows them to tackle complex, multi-file software modifications. Benchmarks such as SWE-bench evaluate agents on their ability to resolve real GitHub issues.

However, almost all these systems operate in a transcript-only paradigm. They focus on improving the success rate of code generation rather than the safety, isolation, or auditability of the execution process. Our work addresses this gap, focusing on the control plane of the agent's environment rather than the model's internal prompt parameters.

\subsection{Program Verification and Sandboxing}
Our custody and proof primitives build on a long history of program verification and systems sandboxing. Traditional verification systems focus on mathematical proof of correctness (e.g., using Coq, Isabelle, or TLA+). In contrast, our proof primitive focuses on pragmatic, machine-checkable verification (compilers, test suites, static analysis) that can be easily executed in standard software engineering workspaces.

Sandboxing technologies (e.g., Docker, WebAssembly, gVisor) are widely used to secure untrusted code execution. In our framework, we apply these concepts to agent tool execution. Decapod's custody primitive uses Git worktrees and shell sandboxing to isolate agent writes, preventing authority leakage and concurrent collisions.

\subsection{Distributed Systems and Ledger Audits}
The framing of agentic coding as a distributed systems execution problem is inspired by classic distributed consensus and audit architectures. Lamport's happenings-before relation~\cite{lamport1978time} laid the groundwork for ordering events across independent computing nodes. MapReduce~\cite{dean2004mapreduce} and distributed databases demonstrated how to partition state and manage concurrency across large scale clusters.

Similarly, cryptographically bound ledgers (such as Git's tree hashes or blockchain transactions) provide inspiration for Decapod's Trajectory and Proof primitives. By logging every tool call and workspace modification in an append-only ledger and signing the result, we bring distributed auditability to the local development environment.
